\chapter{おわりに}\label{cha:Conclusion}
% 本研究では、テストケースの可視化を目的として、
% テストケース可視化ツール\tool{}(VIsualization of Testcase Tool)を試作した。
% \tool{}は、要求仕様書とテストスイートを入力として受け取り、
% 本研究で提案した、TCVダイアグラムを出力する。
% TCVダイアグラムは、テストケースの可視化を目的としたダイアグラムであり、以下の4つを可視化する。

% \begin{itemize}
%     \item テストケースの充足:ある要求仕様を検証するテストケースが存在する状態
%     \item テストケースの抜け:ある要求仕様を検証するテストケースが存在しない状態
%     \item テストケースの不備:ある要求仕様を検証するテストケースは存在するが、内容が要求仕様に対して不十分な状態
%     \item テストケースの過剰:テストケースは存在するが、対応する要求仕様が存在しない状態
% \end{itemize}

% 本研究で試作した\tool{}に、「閏年判定システム」を対象とした適用例を用いて、
% 要求仕様書とテストスイートを入力することで、
% \tool{}が正しくTCVダイアグラムを出力することを確認した。

% \tool{}の有用性を確認するため、\tool{}を使用した場合と使用しなかった場合で評価実験を行った。
% 評価実験では、意図的に含めたテストケースの抜けや不備または過剰箇所を、
% 被験者が全て発見するまでにかかる時間を測定した。
% この結果、組$\alpha$（閏年判定システム）に対して、
% テストケースの抜けや不備または過剰箇所の発見にかかる時間を4分28秒（81.5\%）削減できた。
% また、組$\beta$（成績付けシステム）に対して、
% テストケースの抜けや不備または過剰箇所の発見にかかる時間を3分54秒（80.7\%）削減できた。
% また、評価実験において、\tool{}を使用した場合、テストケースの抜けや不備または過剰箇所の発見漏れはなかったが、
% \tool{}を使用しなかった場合、テストケースの抜けや不備または過剰箇所の発見漏れがあった。
% このことから、\tool{}が出力するTCVダイアグラムは、
% テストケースの抜けや不備または過剰箇所の発見に有用であることを確認した。

% 以上のことから、\tool{}は、テストケースの可視化に有用であるといえる。

% 以下に、今後の課題を示す。

